\lysh{34146}{2}

\begin{alist}
\item 
\begin{rlist}
\item Για $\alpha = 1$ η εξίσωση γίνεται:
\[(1 + 3)x = 1^2 - 9 \Leftrightarrow 4x = -8 \Leftrightarrow x = -2.\]

\item Για $\alpha = -3$ η εξίσωση γίνεται:
$(-3 + 3)x = (-3)^2 - 9 \Leftrightarrow 0x = 0$, που έχει άπειρες λύσεις.
\end{rlist}

\item Η εξίσωση έχει μοναδική λύση αν και μόνο αν: $\alpha + 3 \ne 0 \Leftrightarrow \alpha \ne -3$.
Για την εύρεση της μοναδικής λύσης της εξίσωσης έχουμε:
\begin{align*}
(\alpha + 3)x = \alpha^2 - 9 \Leftrightarrow (\alpha + 3)x = (\alpha + 3)(\alpha - 3) \ &\Leftrightarrow \\
x &= \dfrac{(\alpha + 3)(\alpha - 3)}{\alpha + 3} \Leftrightarrow x = \alpha - 3.
\end{align*}
\end{alist}
